% Coefficients d'étalonnage k_f et k_m (couple positif) en fonction de la position (x, y)
% avec le plan k = 1 (étalonnage constructeur) en référence.
% Données : images/data/Kf.csv et images/data/Km_positif.csv
%   3 colonnes x (mm), y (mm), k, séparées par des virgules, sans en-tête,
%   x variant le plus vite (11 valeurs de x par ligne de la grille), NaN pour une vis sans valeur.
% Nécessite \usepgfplotslibrary{groupplots}.
\begin{tikzpicture}
\pgfplotsset{
    surfk/.style={surf, shader=faceted interp, faceted color=black!25, line width=0.1pt,
                  mesh/cols=11, unbounded coords=jump, opacity=0.95},
    planun/.style={surf, shader=flat, draw=jvcorail, line width=0.4pt,
                   fill=jvcorail, fill opacity=0.18, samples=2, domain=-150:150, y domain=-150:150},
}
\begin{groupplot}[
    group style={group size=2 by 1, horizontal sep=16mm},
    width=0.46\linewidth, height=0.40\linewidth,
    view={-35}{25},
    xmin=-150, xmax=150, ymin=-150, ymax=150,
    zmin=0.75, zmax=1.1,
    xtick={-100,0,100}, ytick={-100,0,100}, ztick={0.8,0.9,1,1.1},
    /pgf/number format/use comma, /pgf/number format/1000 sep={},
    every axis/.append style={font=\scriptsize},
    tick label style={font=\tiny},
    label style={font=\scriptsize},
    title style={font=\scriptsize\bfseries, text=jvbleu},
    xlabel={$x$ (mm)}, ylabel={$y$ (mm)},
    grid=major, grid style={jvgris!20},
    colormap/viridis, point meta min=0.75, point meta max=1.1,
]
    \nextgroupplot[title={Force : $k_f$}, zlabel={$k$}]
    \addplot3[surfk] table[header=false, col sep=comma, x index=0, y index=1, z index=2]
        {images/data/Kf.csv};
    \addplot3[planun] {1};

    \nextgroupplot[title={Couple : $k_m$},
        colorbar, colorbar style={width=2mm, ymin=0.75, ymax=1.1, ytick={0.8,0.9,1,1.1},
                                   yticklabel style={font=\tiny}}]
    \addplot3[surfk] table[header=false, col sep=comma, x index=0, y index=1, z index=2]
        {images/data/Km_positif.csv};
    \addplot3[planun] {1};
\end{groupplot}
% légende du plan
\node[anchor=north, font=\tiny, text=jvcorail] at ([yshift=-9mm]$(group c1r1.south)!0.5!(group c2r1.south)$)
    {\tikz\fill[jvcorail, fill opacity=0.3, draw=jvcorail] (0,0) rectangle (3mm,2mm);~plan $k = 1$};
\end{tikzpicture}
